%------------------------------------------------------------------------------%
\documentclass[a4paper,12pt,fleqn]{article}

\usepackage{calculo_varias_variaveis-1}

\ShowAnswers

%------------------------------------------------------------------------------%
\begin{document}

\makeHeader{CFVVI}{Prova 1 -- Turma 3}{01/04/2026}

% 01
%------------------------------------------------------------------------------%
\begin{multicols}{2}
\raggedbottom
\exercise{50\%}
  Uma partícula descreve uma trajetória circular descrita
  pelas equações paramétricas
  \begin{gather*}
    x(t) = 3 \sin\left(t^2\right) +3 \\
    y(t) = 3 \cos\left(t^2\right) +1 \\
    0 \leq t \leq \sqrt{\pi}
  \end{gather*}
  \begin{enumerate}[label=\alph*),itemsep=0pt,topsep=0pt]
    \item\!\! [10\%] determine on pontos onde da curva associados com \\
                     \hspace*{29pt}
                     $t=0$,
                     $t=\sqrt{\frac{\pi}{2}}$ e
                     $t=\sqrt{\pi}$
    \item\!\! [15\%] calcule a expressão para o vetor tangente à curva
    \item\!\! [15\%] avalie o vetor tangente em \\
                     \hspace*{29pt}
                     $t=\sqrt{\frac{\pi}{2}}$ e
                     $t=\sqrt{\pi}$
    \item\!\! [10\%] esboce a curva e os vetores
                     tangentes calculados
  \end{enumerate}
\columnbreak
~

\begin{questiononly}
  \hfill\includegraphics{axis_xy}
\end{questiononly}
\begin{answeronly}
  \hfill\includegraphics{T3_B_parametrica}
\end{answeronly}
\end{multicols}
\textit{
  Apenas para esboçar os vetores, utilize as aproximações grosseiras
  $\sqrt{\frac{\pi}{2}} \approx 1$ e
  $\sqrt{\pi}           \approx \frac{3}{2}$
}
\clearpagequestiononly

\begin{answer}
  \textbf{a)} \quad
  Calculando os pontos associados a $t=0$, $t=\sqrt{\frac{\pi}{2}}$ e em $t=\sqrt{\pi}$
  \begin{multicols}{2}
    \begin{align*}
      x(0)
      & = \left(3 \sin\left(t^2\right) +3\right) \evalat{0}{} \\
      & = 3 \sin\left(0\right) +3 \\
      & = 3 \times 0 +3 \\
      & = 3
    \end{align*}
    \newcolumn
    \begin{align*}
      y(0)
      & = \left(3 \cos\left(t^2\right) +1\right) \evalat{0}{} \\
      & = 3 \cos\left(0\right) +1 \\
      & = 3 \times 1 +1 \\
      & = 4
    \end{align*}
  \end{multicols}

  \clearpage
  \begin{multicols}{2}
    \begin{align*}
      x\left(\sqrt{\frac{\pi}{2}}\right)
      & = \left(3 \sin\left(t^2\right) +3\right) \evalat{\sqrt{\frac{\pi}{2}}}{} \\
      & = 3 \sin\left(\frac{\pi}{2}\right) +3 \\
      & = 3 \times 1 +3 \\
      & = 6
    \end{align*}
    \newcolumn
    \begin{align*}
      y\left(\sqrt{\frac{\pi}{2}}\right)
      & = \left(3 \cos\left(t^2\right) +1\right) \evalat{\sqrt{\frac{\pi}{2}}}{} \\
      & = 3 \cos\left(\frac{\pi}{2}\right) +1 \\
      & = 3 \times 0 +1 \\
      & = 1
    \end{align*}
  \end{multicols}
  \vspace{-\baselineskip}
  \begin{multicols}{2}
    \begin{align*}
      x\left(\sqrt{\pi}\right)
      & = \left(3 \sin\left(t^2\right) +3\right) \evalat{\sqrt{\pi}}{} \\
      & = 3 \sin\left(\pi\right) +3 \\
      & = 3 \times 0 +3 \\
      & = 3
    \end{align*}
    \newcolumn
    \begin{align*}
      y\left(\sqrt{\pi}\right)
      & = \left(3 \cos\left(t^2\right) +1\right) \evalat{\sqrt{\pi}}{} \\
      & = 3 \cos\left(\pi\right) +1 \\
      & = 3 (-1) +1 \\
      & = -2
    \end{align*}
  \end{multicols}
  \vspace{-2\baselineskip}
  \textbf{b)} \quad
  O vetor tangente é
  \(
  v(t) = \Vetor{\frac{dx}{dt}}{\frac{dy}{dt}}
  \)

  Avaliando cada derivada
  \begin{multicols}{2}
    \begin{align*}
      \frac{dx}{dt}
      & = \frac{d}{dt}\left[3 \sin\left(t^2\right) +3\right]   \\[2mm]
      & = 3 \frac{d}{dt}\left[\sin\left(t^2\right)\right]      \\[2mm]
      & = 3 \cos\left(t^2\right) \frac{d}{dt}\left(t^2\right)  \\[2mm]
      & = 3 \cos\left(t^2\right) 2t                            \\[2mm]
      & = 6 t \cos\left(t^2\right)
    \end{align*}
    \begin{align*}
      \frac{dy}{dt}
      & = \frac{d}{dt}\left[3 \cos\left(t^2\right) +1\right] \\[2mm]
      & = 3\frac{d}{dt}\left[\cos\left(t^2\right)\right]     \\[2mm]
      & = -3\sin\left(t^2\right)\frac{d}{dt}\left(t^2\right) \\[2mm]
      & = -3\sin\left(t^2\right)2t \\[2mm]
      & = -6t \sin\left(t^2\right)
    \end{align*}
  \end{multicols}
  concluímos que,
  \(
  v(t) = \vetor%
  { 6t \cos\left(t^2\right)}
  {-6t \sin\left(t^2\right)}
  \)

  \bigskip
  \textbf{c)} \quad
  Avaliando o vetor tangente nos pontos
  $t=\sqrt{\frac{\pi}{2}}$ e $t=\sqrt{\pi}$
  \begin{multicols}{2}
    \begin{align*}
      v\left(\sqrt{\frac{\pi}{2}}\right)
      & = \vetor
      { 6t \cos\left(t^2\right)}
      {-6t \sin\left(t^2\right)}
      \evalat{\sqrt{\frac{\pi}{2}}}{} \\
      & = \vetor
      { 6\sqrt{\frac{\pi}{2}} \cos\left(\frac{\pi}{2}\right)}
      {-6\sqrt{\frac{\pi}{2}} \sin\left(\frac{\pi}{2}\right)}
      \\
      & = \vetor
      { 6\sqrt{\frac{\pi}{2}} \times 0}
      {-6\sqrt{\frac{\pi}{2}} \times 1}
      \\
      & = \vetor{0}{-6\sqrt{\frac{\pi}{2}}}
    \end{align*}
    \newcolumn
    \begin{align*}
      v\left(\sqrt{\pi}\right)
      & = \vetor
      { 6t \cos\left(t^2\right)}
      {-6t \sin\left(t^2\right)}
      \evalat{\sqrt{\pi}}{} \\
      & = \vetor
      { 6\sqrt{\pi} \cos\left(\pi\right)}
      {-6\sqrt{\pi} \sin\left(\pi\right)}
      \\
      & = \vetor{-6\sqrt{\pi}}{0}
    \end{align*}
  \end{multicols}
\end{answer}

% 02
%------------------------------------------------------------------------------%
\exercise{50\%}
Considerando a função
\(
f(x, y) = \ln\left(9+x^2-y^2\right)
\)
\begin{enumerate}[label=(\alph*),itemsep=0pt]
  \item \ [15\%] determine o domínio de $f$
  \item \ [15\%] esboce o domínio
  \item \ [10\%] determine a curva de nível que passa pelo ponto \((0, 1)\)
  \item \ [10\%] esboce a curva de nível que passa pelo ponto \((0, 1)\)
\end{enumerate}
\begin{questiononly}
  \begin{multicols}{2}
    \;\; Domínio
    \vspace{-0.3\baselineskip}
    \begin{center}
      \includegraphics{axis_xy}
    \end{center}
    \;\; Curva de nível
    \vspace{-0.3\baselineskip}
    \begin{center}
      \includegraphics{axis_xy}
    \end{center}
  \end{multicols}
\end{questiononly}
\begin{answeronly}
  \begin{multicols}{2}
    \;\; Domínio
    \vspace{-0.3\baselineskip}
    \begin{center}
      \includegraphics{T3_B_dominio}
    \end{center}
    \;\; Curva de nível
    \vspace{-0.3\baselineskip}
    \begin{center}
      \includegraphics{T3_B_curva_nivel}
    \end{center}
  \end{multicols}
\end{answeronly}

\begin{answer}
  \textbf{a)} \quad
  Não podemos calcular logaritmo de um número negativo ou do zero,
  portanto, o domínio de $f$ é o conjunto
  de todos os pontos \((x, y)\in\R^2\) tais que
  \begin{align*}
    9+x^2-y^2                         & > 0 \\
    -x^2 + y^2                        & < 9 \\
    \frac{y^2}{3^2} - \frac{x^2}{3^2} & < 1
  \end{align*}
  ou seja, os pontos ``entre os lados'' da hipérbole.
  Assim, o domínio é
  \[
  D_f = \{(x, y) \in \R^2 \st y^2 - x^2 < 9\}
  \]

  \bigskip
  \textbf{c)} \quad
  A curva de nível que passa pelo ponto $(0, 1)$ tem o valor
  \[
    c
    = f(0, 1)
    = \ln\left(9 +x^2 -y^2 \right) \evalat{(0, 1)}{}
    = \ln\left(9 +0^2 -1^2 \right)
    = \ln 8
  \]
  A curva de nível associada é
  \begin{align*}
    f(x, y)                   & = \ln 8     \\[2mm]
    \ln\left(9+x^2-y^2\right) & = \ln 8     \\[2mm]
    9 + x^2 - y^2             & = e^{\ln 8} \\[2mm]
    x^2 - y^2                 & = 8 - 9     \\[2mm]
    y^2 - x^2                 & = 1
  \end{align*}
  Que é uma hipérbole.
\end{answer}

%------------------------------------------------------------------------------%
\end{document}
%------------------------------------------------------------------------------%
